\section{March 16}

\textbf{continue from last time}

\subsection{Some remarks of local reciprocity law}

  \begin{remark}
    If $\characteristic K = p > 0$, then there always exists a subgroup of $K^\times$ which is not open.

    Consider the derivation $\delta: K=\bbF_p((t)) \to K$ by $t \mapsto 1$. 
    Let $L/K$ be a finite separable extension, then $\delta$ extends to a derivation $\delta_L: L \to L$.
    (one can view a derivation as a vector field, and the extension is an \'etale covering, since it is locally isomorphic, the vector field can be lifted to the covering.)
    By considering the kernel of $\delta$, we get $L^p \cap K = K^p$.
    Hence $K^\times / (K^\times)^p \subset K^\times / (L^\times)^p$ are $\bbF_p$-vector spaces.
    By direct computation, we have
    \[ K^\times / (K^\times)^p \cong \Image \delta = \left\{ \sum_{i \geq 1} a_i t^i : a_i \in \bbF_p, p \nmid i+1 \right\}. \]
    As $\bbF_p$-vector spaces, $\Image \delta$ has uncountable dimension.
    On the other hand, there exists countable and dense subset $S \subset L^\times$.
    Its image in $L^\times / (L^\times)^p$ is also countable. 
    We can find a codimension one subspace $V \subset L^\times / (L^\times)^p$ containing the image of $S$.
    Then the preimage of $V$ in $L^\times$ is a subgroup of $L^\times$ of index $p$ which is dense in $L^\times$.
    It is not closed and hence not open by finite index.
  \end{remark}

  \begin{remark}\label{rmk:unit_has_trivial_image_under_artin_local_map_when_restrict_to_unramified_extension}
    Let 
    \[ K^\unr \coloneqq \bigcup_{L/K \text{ finite unramified}} L. \]
    The residue field of $K^\unr$ is an algebraic closure of the residue field of $K$ and $\Gal(K^\unr/K) \cong \Gal(\overline{k}/k)$.
    Then there exists a unique element $\Frob_K \in \Gal(K^\unr/K)$ such that $\Frob_K$ induces the $q$-th power map on the residue field, where $q = \# k$.
    We have an isomorphism of topological groups
    \[ \hat{\bbZ} \to \Gal(K^\unr/K), \quad \alpha \mapsto \Frob_K^\alpha. \]
    Then the (a) in \cref{thm:local_reciprocity_law} can be restated as follows: for all uniformizer $\pi$ of $K$, $\phi_K(\pi)|_{K^\unr} = \Frob_K$.
    In particular, for all unit $u \in \calO_K^\times$, $\phi_K(u)|_{K^\unr} = \id$.
  \end{remark}

  \begin{remark}
    The groups 
    \[ U_K \supset 1 + \frakm_K \supset 1 + \frakm_K^2 \supset \cdots \]
    form a fundamental system of neighborhoods of $1$ in $U_K$.
    Let $L/K$ be a finite abelian extension.
    If $L/K$ is unramified, then $U_K \subset \ker (\phi_{L/K})$ by \cref{rmk:unit_has_trivial_image_under_artin_local_map_when_restrict_to_unramified_extension}.
    We say that $L/K$ has \emph{conductor} $0$.
    Otherwise, the smallest integer $n \geq 1$ such that $1 + \frakm_K^n \subset \ker (\phi_{L/K})$ is called the \emph{conductor} of $L/K$.
    
    Tautologically, the extension $L/K$ is unramified if and only if its conductor is $0$.
    The conductor measures the ramification of $L/K$.
    We say that $L/K$ is \emph{tamely ramified} if its conductor is $\leq 1$, and \emph{wildly ramified} otherwise.
    Note that $(1+\frakm^i)/(1+\frakm^{i+1})$ are $p$-groups for $i \geq 1$, so $\phi|_{L/K}|_{U_K}$ factors through $U_K/(1+\frakm_K)$ iff $L/K$ is tamely ramified.
  \end{remark}

  \begin{remark}
    The homomorphisms
    \[ \phi_{L/K}: K^\times / \Nm_{L/K}(L^\times) \to \Gal(L/K) \]
    form an inverse system as $L$ varies over all finite abelian extensions of $K$, ordered by inclusion.
    By taking the inverse limit, we get an isomorphism of topological groups
    \[ \widehat{\phi}_K: \widehat{K^\times} \to \Gal(K^\ab/K), \]
    where $\widehat{K^\times}$ is the profinite completion of $K^\times$.
    This topology on $K^\times$ is called the \emph{norm topology}.
    
    By \cref{thm:norm_groups_are_exactly_open_subgroups_of_finite_index}, norm topology is coarser than the original topology on $K^\times$.
    Fix a uniformizer $\pi$ of $K$, then $K^\times \cong \pi^\bbZ \cdot U_K$.
    Then subgroups $(1+\frakm^n)\left<\pi^m\right> \cong (1+\frakm^n) \times m\bbZ$ form a fundamental system of neighborhoods of $1$ in the norm topology.
    In particular, $U_K$ is not open under the norm topology.
    When completing $K^\times$ under the norm topology, we get $\widehat{K^\times} = \pi^{\hat{\bbZ}}\cdot U_K \cong \hat{\bbZ} \times U_K$.

    More canonically, the identity map 
    \[ (K^\times, \text{ original topology}) \to (K^\times, \text{ norm topology}) \] 
    is continuous and induces a homomorphism of completion $K^\times \to \widehat{K^\times}$ fitting into the following commutative diagram
    \[ \begin{tikzcd}
        0 \arrow[r] & U_K \arrow[r] \arrow[d, equal] & K^\times \arrow[r] \arrow[d] & \bbZ \arrow[r] \arrow[d, hook] & 0 \\
        0 \arrow[r] & U_K \arrow[r] & \widehat{K^\times} \arrow[r] & \hat{\bbZ} \arrow[r] & 0.
    \end{tikzcd} \]
  \end{remark}

  \begin{remark}\label{rmk:decomposition_of_K_ab_into_unramified_and_totally_ramified_part}
    The choice of a uniformizer $\pi$ of $K$ gives a decomposition $\widehat{K^\times} \cong \hat{\bbZ} \times U_K$ of $\widehat{K^\times}$ into a product of two closed subgroups.
    Hence we get a decomposition of $K^\ab = K^\unr \cdot K_\pi$, where $K_\pi$ is the subfield of $K^\ab$ fixed by $\phi_K(\pi)$ and $K^\unr$ is the subfield of $K^\ab$ fixed by $\phi_K(U_K)$.
    Clearly, 
    \[ K_\pi = \bigcup_{L/K \text{ finite abelian}, \pi \in \Nm(L^\times)} L. \]
    Here the condition $\pi \in \Nm(L^\times)$ implies that $L/K$ is totally ramified.
    The map $\phi_K$ restricts to an isomorphism $U_K \to \Gal(K_\pi/K)$.
    It maps the filtration $U_K \supset 1 + \frakm_K \supset 1 + \frakm_K^2 \supset \cdots$ to the filtration of $\Gal(K_\pi/K)$ by ramification groups.
  \end{remark}

  \begin{example}
    When $K = \bbQ_p$ with uniformizer $p$, the decomposition in \cref{rmk:decomposition_of_K_ab_into_unramified_and_totally_ramified_part} is given by $\bbQ_p^\ab = \bbQ_p^\unr \cdot \bbQ_p(\mu_{p^\infty})$, where $\bbQ_p(\mu_{p^\infty})$ is the union of all cyclotomic extensions of $\bbQ_p$ and $\bbQ_p^\unr = \bigcup_{p \neq m} \bbQ_p(\zeta_m)$.
  \end{example}

  \begin{remark}
    The composition of two totally ramified extensions need not be totally ramified.
    For example, pick $u \in \bbZ_p^\times$ which is not a square.
    Both $\bbQ_p(\sqrt{p})/\bbQ_p$ and $\bbQ_p(\sqrt{pu})/\bbQ_p$ are totally ramified, but $\bbQ_p(\sqrt{p}, \sqrt{pu})/\bbQ_p$ is not totally ramified since it contains $\bbQ_p(\sqrt{u})/\bbQ_p$ which is unramified.
  \end{remark}


\subsection{Outline of the proofs of the main theorems in local CFT}

  \begin{theorem}[uniqueness of the local reciprocity map]\label{thm:uniqueness_of_local_reciprocity_map}
    Assume \cref{thm:norm_groups_are_exactly_open_subgroups_of_finite_index}, then there exists at most one homomorphism $\phi_K: K^\times \to \Gal(K^\ab/K)$ satisfying the properties in \cref{thm:local_reciprocity_law}.
  \end{theorem}
  \begin{proof}
    Assume there exists such homomorphism $\phi_K$.
    The \cref{cor:correspondence_between_extensions_and_norm_groups} holds with ``norm groups'' replaced by ``open subgroups of finite index''.
    Let $\pi$ be a uniformizer of $K$, $K_{\pi,n}$ be the extension of $K$ with norm group $(1+\frakm_K^n)\left<\pi\right>$, and $K_\pi = \bigcup_{n \geq 1} K_{\pi,n}$.

    As $\pi$ is a norm, $\phi_K(\pi)$ acts trivially on $K_\pi$.
    Since it acts as the Frobenius on $K^\unr$, the action of $\phi_K(\pi)$ on $K_\pi K^\unr$ is completely determined.
    On the other hand, $K_\pi K^\unr = K^\ab$ by \cref{rmk:decomposition_of_K_ab_into_unramified_and_totally_ramified_part}.
    Hence $\phi_K(\pi)$ is determined by the properties in \cref{thm:local_reciprocity_law}, 
    i.e., $\phi_K(\pi) = \phi'_K(\pi)$ for any other homomorphism $\phi'_K$ satisfying the properties in \cref{thm:local_reciprocity_law}.
    Note that we can change the uniformizer $\pi$ by any element in $\pi U_K$, and the same argument shows that $\phi_K(\pi u) = \phi'_K(\pi u)$ for all $u \in U_K$.
    Hence $\phi_K = \phi'_K$.
  \end{proof}

  Hence to prove the local reciprocity law, it suffices to show the existence in \cref{thm:local_reciprocity_law} and to prove \cref{thm:norm_groups_are_exactly_open_subgroups_of_finite_index}.

  There are three proofs.
  The first two give an explicit construction of $K_{\pi,n}$ and the local Artin map $\phi_{K_{\pi,n}/K}$.
  The last two will introduce some cohomological results and these results will be used in the global CFT as well.
  The last one is completely cohomological and does not give an explicit construction of the local Artin map, but it is simpler.

  The first proof is based on Lubin-Tate formal groups. It gives an explicit construction of $K_{\pi,n}$ and the local Artin map 
  \[ \phi_{K_{\pi}/K}: U_K \to \Gal(K_\pi/K). \]
  Moreover, it shows that $K_\pi \cdot K^\unr$ and $\phi|_{K_\pi \cdot K^\unr}$ are independent of the choice of $\pi$.
  Finally, it uses the Hasse-Arf theorem to show that $K_\pi \cdot K^\unr = K^\ab$.